\section{Method}

\subsection{Overview}
\textsc{Struct Memory-R1} instantiates the memory-augmented LLM framework defined in Section~2 with three learned components: a Memory Manager $\pi_{\theta_M}$, a Retrieve Agent $\pi_{\theta_R}$, and an Answer Agent $\pi_{\theta_A}$. Each component is parameterized by a language model and trained with Group Relative Policy Optimization (GRPO)~\cg{shao2024deepseekmath} using downstream question-answering reward. Training proceeds in stages: the Answer Agent is trained first, then frozen while the Retrieve Agent and Memory Manager are trained in sequence. This staged approach ensures stable reward signals at each phase.

\subsection{Memory Bank and Tree Structure}
The memory bank $\mathcal{T}_t = (V_t, E_t)$ is a rooted tree that is constructed incrementally from raw conversation. Each node $v \in V_t$ resides at a level $\ell(v) \in \{0, 1, \ldots, L\}$ and carries content $c(v)$. The tree is structured such that:
\begin{itemize}[nosep,leftmargin=1.5em]
\item Level 0: the root, representing the full conversation;
\item Levels $1$ to $L{-}1$: internal nodes representing hierarchical groupings (e.g., people, topics, events);
\item Level $L$: leaf nodes, each containing an individual memory entry $m_i$.
\end{itemize}
Each internal node $v$ maintains a summary $s(v)$ that aggregates information from its children $\text{Ch}(v) = \{u \in V_t : (v, u) \in E_t\}$. The \emph{schema} of the tree, $\textsc{Schema}(\mathcal{T}_t)$, consists of node identifiers and summaries at each level without exposing full leaf content, providing a compact representation for the retrieval policy.

The tree is serialized as structured JSON for input to language model policies. This serialization preserves the hierarchical relationships while remaining compatible with standard text-based LLM input formats.

\subsection{Memory Manager}
The Memory Manager observes the current dialogue turn $x_t$ and the current memory tree $\mathcal{T}_t$, and emits a set of operations to update the tree. We restrict the operation space to four stable CRUD operations:
\begin{equation}
\mathcal{O} = \{\textsc{Add}, \textsc{Update}, \textsc{Delete}, \textsc{None}\}.
\end{equation}
Given turn $x_t$ and tree $\mathcal{T}_t$, the manager samples an operation sequence:
\begin{equation}
o_t \sim \pi_{\theta_M}(\cdot \mid x_t, \mathcal{T}_t),
\end{equation}
and the updated tree is obtained deterministically:
\begin{equation}
\mathcal{T}_{t+1} = \textsc{Apply}(\mathcal{T}_t, o_t).
\end{equation}
The operation $o_t$ is a JSON object specifying which entries to add, modify, or remove, along with their placement in the tree hierarchy. By restricting to CRUD operations rather than arbitrary structural rewrites, we isolate the benefit of structured memory from the complexity of unrestricted topology editing.

\paragraph{Reward.}
The Memory Manager is trained with outcome-driven reward. After applying the manager's operations, the downstream pipeline (Retrieve Agent + frozen Answer Agent) processes the corresponding question. The reward is:
\begin{equation}
r_M = \mathcal{R}(y_t, y_t^\star) + \lambda_f \cdot \mathbb{1}[\text{valid JSON}],
\end{equation}
where $\mathcal{R}(y_t, y_t^\star)$ measures answer quality and $\lambda_f$ is a small bonus for well-formed output. This encourages the manager to write memories that are useful for future question answering rather than those matching local heuristics.

\subsection{Retrieve Agent}
The Retrieve Agent operates over the memory tree constructed by the Memory Manager. Given a question $q$ and the tree schema $\textsc{Schema}(\mathcal{T})$, it emits a search plan:
\begin{equation}
z \sim \pi_{\theta_R}(\cdot \mid q, \textsc{Schema}(\mathcal{T})).
\end{equation}
The search plan $z$ is a JSON object of the form:
\begin{equation}
z = \{\texttt{levels}: [\ell_1, \ldots, \ell_k], \;\; \texttt{selected\_ids}: [id_1, \ldots, id_m], \;\; \texttt{stop}: \{0,1\}\},
\end{equation}
where \texttt{levels} specifies which branches to expand at each tree level, \texttt{selected\_ids} names specific leaf nodes to retrieve, and \texttt{stop} terminates search. A deterministic executor runs the plan against the tree:
\begin{equation}
R = \textsc{Execute}(\mathcal{T}, z),
\end{equation}
returning a set of retrieved memory leaves together with local context (parent summaries, sibling information). This design ensures that retrieval is conditioned on structural locality rather than global similarity scoring.

\paragraph{Reward.}
The retriever reward combines answer utility with evidence quality:
\begin{equation}
r_R = \mathcal{R}(y_t, y_t^\star) + \lambda_{\text{ev}} \cdot \textsc{Hit}(R, E^\star) - \lambda_{\text{cost}} \cdot C(z),
\end{equation}
where $E^\star$ denotes gold evidence nodes (when available), $C(z)$ penalizes search cost, and the primary signal is answer quality under a frozen Answer Agent.

\subsection{Answer Agent}
The Answer Agent receives the question $q$ and retrieved memories $R$, and produces a response:
\begin{equation}
y \sim \pi_{\theta_A}(\cdot \mid q, R).
\end{equation}
This component is trained to condition on retrieved evidence rather than on the full conversation history. During Memory Manager and Retrieve Agent training, the Answer Agent is frozen and serves as the evaluator in the reward loop, ensuring that improvements in upstream components are measured through their effect on answer quality.

\subsection{Training with GRPO}
All three components are trained using Group Relative Policy Optimization (GRPO)~\cg{shao2024deepseekmath}, which avoids the need for a separate value network by computing advantages from a group of sampled outputs. For a given input $x$ and policy $\pi_\theta$, GRPO samples a group of $G$ outputs $\{o_1, \ldots, o_G\}$ and computes normalized advantages:
\begin{equation}
\hat{A}_i = \frac{r_i - \mu(\{r_1, \ldots, r_G\})}{\sigma(\{r_1, \ldots, r_G\})},
\end{equation}
where $r_i$ is the reward for output $o_i$, and $\mu, \sigma$ denote the group mean and standard deviation. The policy is updated by maximizing:
\begin{equation}
\mathcal{L}_{\text{GRPO}}(\theta) = \mathbb{E}_{x} \left[ \frac{1}{G} \sum_{i=1}^{G} \min\!\left( \rho_i \hat{A}_i, \; \text{clip}(\rho_i, 1{-}\epsilon, 1{+}\epsilon) \hat{A}_i \right) - \beta \, D_{\text{KL}}\!\left[\pi_\theta \| \pi_{\text{ref}}\right] \right],
\end{equation}
where $\rho_i = \pi_\theta(o_i \mid x) / \pi_{\text{old}}(o_i \mid x)$ is the importance ratio, $\epsilon$ is the clipping parameter, and $\beta$ controls the KL penalty relative to a reference policy $\pi_{\text{ref}}$. This objective encourages the model to produce outputs that yield higher rewards relative to the group average while staying close to the reference distribution.
